\chapter[Sermon 52. The Description of the Church and of the Red Dragon, Fighting Against the Church.]{The Description of the Church and of the Red Dragon, Fighting Against the Church.}
\label{sermon:052}
\markright{Sermon 52. The Description of the Church and of the Red Dragon ...}

And there appeared a great sign in the heavens: a woman clothed with the sun, and the moon under her feet, and upon her head a crown of twelve stars. And she was with child, and cried out in labor, in pain ready to give birth. And another sign appeared in the heavens: behold, a great red dragon, having seven heads, and ten horns, and seven crowns upon his heads. And his tail drew a third of the stars of heaven and cast them to the earth. And the dragon stood before the woman, who was ready to give birth, to devour her child as soon as it was born. And she brought forth a male child, who would rule all nations with a rod of iron. And her son was taken up to God and to his throne. And the woman fled into the wilderness, where she had a place prepared by God, that they should nourish her there for one thousand two hundred and sixty days.

\index{Antichrist}The fourth section of this book presents to us the third vision, which others who divide the second vision into two consider to be the fourth. The Lord has frequently and extensively mentioned in the second vision the persecution and conflict of the faithful with Antichrist and with God’s wicked enemies, especially in chapters 6, 9, and 11. He now proceeds in the third vision, and at length, to discuss this conflict, setting forth the events as if they were to be seen with the eyes, in the three chapters immediately following chapters 12, 13, and 14. He repeats all things more deeply, and vividly and diligently describes the aspects of this conflict, and afterward the fight itself. Therefore, after describing the church, which endures the brunt of this war, he also describes the Dragon that instigates the war. He declares how busily he watches, and again, lest anyone should be discouraged, he adds that, despite his efforts, Christ certainly overcomes him, and that God ultimately hinders and defeats his endeavors, yielding him vanquished to the faithful. Now he describes the chief instruments that Satan uses in assaulting and persecuting the church, namely the old and new Roman Empire, and within this, the corrupt papistry, in which Antichrist is also vividly portrayed. Nevertheless, he adds to these unfortunate events, for the consolation and comfort of the godly, that the Lamb nevertheless stands on Mount Zion as a conqueror, having his church with him, however much this world rages and may be never so mad and cruel, and that the gospel is preached in defiance of Antichrist, and all men are warned to beware of Antichrist. Here he also begins to reason about God’s judgment against the wicked, so as to prepare him to speak in the fifth section about the pains or punishments of the Antichristians, a treatise which he begins in chapter 15. Hitherto, therefore, he treats of the fight or conflict of the church and of the wicked, namely of Antichrist, all of which the father of all murder and of all iniquity, the Devil, inspires.

Therefore, just as this entire book is drawn from scripture and explains the older scriptures very well, these things also which were mentioned earlier seem to be taken from the third chapter of Genesis. There the Lord says, “I will put enmity between you (meaning the Serpent) and the woman, between your seed and her seed: her seed shall bruise your head, and you shall bruise his heel.” For you will read at the end of that chapter as well: “And the dragon was angry with the woman and departed to make war with the remnant of her seed.”

And he describes above all things the parts of this conflict, her who was besieged by war, and the one who instigated the war, namely, the church and the Dragon. And he says how a sign of these things appeared in the heavens. For he would not only say or write, but also set them forth to be seen by the eyes, and in a manner to depict, so that all things might be seen more plainly. And where he says those signs were great, he admonishes that they were and are things of the greatest weight, and matters of the greatest importance.

\index{Primasius}First, he describes the church of God throughout all times using the image or figure of a woman. It is neither strange nor rare that the woman began to represent the image of Christ’s bride, the church, at the very beginnings of things, as can be seen in Genesis 2. And so the Apostle explained the image in chapter 5 of Ephesians. I need not now recount that Isaiah has repeatedly used the image of a woman to represent the church of God: “Rejoice, you barren one who does not bear,” he says, and so on. Finally, Paul in Galatians 4 uses Sarah as a figure of the church, which Solomon also explored at length in his Song of Songs, describing his bride. The church is that woman united with Christ, her spouse, in true faith and continual love. Afterward, he applies certain things specifically to the Virgin Mary, although the things that precede and follow do not entirely agree with this, a point shown by Methodius and Primasius, and other commentators as well.

This woman is clothed with the sun. Scripture calls Christ the sun of righteousness and the light of life. Paul commands the church to put on Christ. Therefore, he is the light, the life, and the righteousness of the church; by Christ, the nakedness of the church is covered; Christ is the ornament and beauty of the church, through him it shines in the world.

The Moon is subject to alterations, is variable, and receives sundry colors: she increases, and decreases: and although it shine, yet it always appears full of spots, and borrows her light of the Sun. Therefore all courses and alterations of times, and whatever is mutable and corruptible in this world, all affections also and infirmities, the church treads under her feet. All the light that she has, she has it of Christ; the light of her righteousness increases and decreases. Finally, she gathers always some spots of the nature of flesh, which she cannot leave but by death. Therefore she shines indeed, however the church feels some obscurity. As the Lord has said also, every branch bearing fruit he purges, that he may bring forth more fruit. And he that is washed, is all clean, and needs no more but to wash his feet.

Moreover, a crown is the honor of the head, and a sign of a kingdom. Christ is the beauty, comeliness, and king of the church. In this crown there are no precious stones, but stars. For in Christ are the patriarchs, prophets, and the twelve apostles, who beautify and enlighten the church, and who have the light of the crown, and dispense that power into the church. Thus is signified the doctrine of the ministers, as in the first chapter of this book. Neither is the shining ministry the smallest portion, but it represents the most excellent things of the church.

Moreover, that woman has in her belly, which in a certain phrase of speech is as much to say as that woman was with child; and had not only a great belly (as we say) but, after the manner of women in labor, cried out and, laboring with pain, was about to be delivered. This properly does not pertain to the Virgin Mary, but to the church. For the primitive church, concerning that first promise of the blessed seed, conceived in her mind a most assured hope that, at length, the Son of God should be born of a virgin, namely, the seed promised, which would crush the serpent’s head. Therefore, the church with an earnest desire and with most fervent prayers coveted and wished that Christ might once be engendered in and by the excellent womb of the same holy virgin. Moreover, Christ is begotten in his faithful when, through his virtue, they are regenerated. For Paul says, “My little children, for whom I labor again, until Christ is shaped in us.” The church therefore travails and brings forth in two ways: bodily, while she earnestly desires, with pain, that Christ might be born of the virgin; and spiritually by faith and regeneration, while she desires to be made conformable to Christ in her members. This therefore is the nature and disposition of this woman, that with a great desire embracing the incarnation of Christ and redemption, she would fain have it known to many; and that many times she wisheth to be regenerated and reformed after the image of Christ.

This is truly a fine description of the church. Compare this to those who today claim the title and the guise of the church, and judge how well they align with this description. But this true church of Christ is brought into danger and conflict.

Now let us hear, in the second place, and as it were on the opposing side, what kind of person is the adversary or enemy of the church: namely, that old serpent, who was a liar and a murderer from the beginning, the sole author of all evil, all mischief, all errors, all iniquity, murder, and disturbance, and a most ungracious devil. Whom afterward he calls Satan, the seducer of the world, and adorns him with other titles, fitting for such a majesty.

This is the Dragon, and that the great Dragon, namely, of great power throughout the world in his followers. And a Dragon, because in ancient times he assumed the form of a serpent and deceived our parents. Pliny and other authors write many things about dragons. Scripture in some places, calls the Devil a twisting serpent. For he is wonderfully subtle and can transform himself into countless forms, that he may deceive and keep the deceived in error.

He is red, for he is full of fire, and the blood of saints and of the innocent. A true bloodhound, the parent and patron of all persecutors and bloody soldiers. In him remain the stains of the blood of Abel. He still smells of the shedding of the blood of the Prophets and Apostles.

This has seven heads; upon each of these is seen a royal crown. He also has ten horns. For the Devil is called the prince of this world and has indeed been the governor of wicked rulers of all ages, and the ringleader of all horns and bloody realms. He was therefore the head of Ninus, the king and prince of Pharaoh, chief captain of Balthazar, king of Babylon, of Cambyses the Persian, of Antiochus the Macedonian, of Julius Caesar the Roman, and likewise of all other tyrants.

The prophet Isaiah called a false prophet a tail, because of his soothing and flattering words; for with his honeyed tongue and sweet words, he creeps into favor with great men. Therefore, with flattering and deceptive words, and lying promises, with which (as in times past) he promises his worshippers godly things, he persuades them to all wickedness, stars—that is to say, preachers and notable men—whom he draws away from heavenly things and casts unto earthly things, so that, having forgotten celestial matters and their holy office and duty, they may now cling to earthly things, being wrapped in the earthly snares of the devil’s tail. And thus in deed he shall corrupt not a few. For he puts the third part of stars, meaning a great number of notable men, whose ministry he uses against the church. Thereof there are many and notable examples from all times in all histories.

And after he has described this foul and filthy beast, and sworn enemy of all saints from the beginning of the world, he also reveals his attempts, treasons, and bitter poison against the church, and how he began to stir up war. This dragon, he says, stood before the woman who was ready to give birth, and he stood watching, diligent, attentive, and eagerly awaiting at all times, and he observed, taking advantage of every opportunity to harm the church. But the end of all his endeavors was to devour the son, born of the spouse of God. He has always, even from the beginning of the world, sought to intercept the glory of Christ; and if any faithful person of the church is spiritually regenerated and made conformed to Christ, he also attempts to lead them into error and destroy them. Therefore, Saint Peter did not say without cause that the Devil goes about like a hungry lion, seeking whom he may devour.

He now, by the way, shows that Christ, as he was promised, is presented to the church, so that the dragon could do nothing against him. Whereupon he wants us to conclude absolutely that he shall have no power over us either, if we abide in Christ. For now he shifts from the universal church to the singular and most excellent member thereof, the virgin Mary, and in few words knits up the mystery of the incarnation: that excellent woman, of whom it is spoken in Genesis 3. The daughter of that matron, I mean the church, the holy Virgin, brought forth a man child, that is to say, her firstborn, king and priest, as Luke testifies in chapter 2. Immediately he declares what and of how great power he is, and why he called him a man child. He is the one of whom David prophesied in Psalm 2, that he should rule all nations with a rod or scepter, not of wood or lead that is pliable, but of iron—that is, strong and durable—namely, the word of God. But those who do not obey God’s word, he will beat down far and never, with an iron staff, that is, with power, which no man is able to resist. But for this mighty prince, Satan laid an ambush, that old dragon, who stirred up against the chief of the Jews and Gentiles. But he found nothing in him at all, as the Lord himself said in John 14; nor shall he find anything in the faithful of Christ at the last. Moreover, while the dragon attempted great things against Christ by the elders of the Jews, being risen from the dead, the Lord was taken up, as it were, out of the throat or hottest assaults of the dragon, to his heavenly Father, and sat on the right hand of God the Father, the old serpent’s attempts made frustrate. And thither also he will receive his faithful, though the serpent’s guts should burst. For through hope we sit together with our head in the supercelestial places, Ephesians 2. And this is the chief and greatest hope of the church in this conflict. For thus he gathers: the dragon most strongly and fiercely invades not only the ancient church but even the very head of the church, the redeemer Christ; yet with his fury, he could nothing prevail. Therefore, he shall no more prevail against his members.

Now he returns again to the church, and says: after the dragon could bring nothing to pass against the son of God, he went and made war against the church, and the church fled into the wilderness. Certainly, Judea in the prophets is compared to a place most frequented; the Gentiles are called a desert or wilderness. Therefore, after Christ’s ascension, the Apostles departing out of Judea, repaired to the Gentiles; yea, and the Jews, inspired by the red dragon, cast out the church out of their limits, which was constrained, as appears in the Acts of the Apostles, to flee unto the Gentiles. And where the Lord has prepared a place for his church, and the church was greatly augmented among the Gentiles, certainly it was through his grace (and by no merit of man). Which prepared the place, which calls, directs, and keeps his sheep, the same has disposed, and yet does dispose for this church ministers or pastors, which may feed it, as the ravens did Elias, all the time that shall be, unto the world’s end. For as for the nobility of those days, I discoursed before. And by this exposition is signified that the dragon shall fight stoutly against the church, so that she shall be compelled to flee; but however much he shall rage against the church, the Lord God shall yet prepare a place on earth, wherein she may dwell safe, and will ever send pastors to feed. He shows moreover, that the flight shall not always be reproachable. The Lord save and keep us. Amen.
